\section{What this framework adds}
\label{sec:scope}

The aim is to make mathematical interfaces reusable inside a weak arithmetic
theory. A successful interface should say which finite data represent an object,
what a construction computes, which facts it needs from its inputs, and why
its proof uses only the permitted induction. Category theory is especially
useful for this purpose: composition, naturality, and universal properties
identify the same proof pattern across algebra, analysis, and finite geometry.
Their use must preserve the computational meaning of the presentations.

The repository inspected for this report is ProveIt, at commit
\begin{center}
\small\texttt{bb86091690bc4f920f046a37cfc87719299b4019}.
\end{center}
Its iPRA directory defines
the target as intuitionistic first-order logic with equality, defining equations
for specified primitive recursive functions, successor axioms, and
quantifier-free induction \cite{repo-main,repo-pr}. The main README catalogs
63 research packets, including twelve foundations packets, and explicitly
distinguishes the kernel-checked ConciseIPRA library from contributed packets
that have not been rerun or verified by the repository intake
\cite{repo-index}. A packet's own build report is evidence of what its authors
claim, not a validation performed in this report.

\subsection{The existing foundation is substantial}

All twelve foundations READMEs were examined, along with selected underlying
proof sections and the actual ConciseIPRA syntax and deduction modules. The
application review is targeted: it includes the Puiseux obstruction and the
Mordell--Weil, Siegel, and Faltings difficulty reports. It is not a verification
of every theorem packet. The source inventory accompanies the archive.

\begin{longtable}{@{}P{.24\linewidth}P{.69\linewidth}@{}}
\toprule
Existing work & Contribution retained and credited\\
\midrule\endhead
Finite contracts and embedding packets & Bounded Booleanization, guarded
operations, quantitative names, finite nets, explicit moduli, and the distinction
between a PR algorithm and a proof of its specification.\\
Finite demand and finite observation packets & Backward precision requests,
changing-parameter recursion, aggregate root data, and finite observations that
avoid an unjustified exact branch choice.\\
Evidence and coherent methods & Finite compatibility problems, coherent
restriction maps, filtered lifting, and explicit finite algebraic invariants.\\
Presentations to witnesses & Finite multiplication identities, rational
enumeration, and complete section lists obtained from appropriate algebraic
certificates.\\
Resources and elimination & Actual PR syntax for some recursions, controlled
descent bounds, finite intersection lists, and auxiliary envelopes whose coverage
remains an independent obligation.\\
\bottomrule
\end{longtable}

The relevant original packets are credited throughout
\cite{repo-demand,repo-observation,repo-evidence,repo-coherent,repo-witnesses,repo-resources}.
Simply repeating that mathematical objects can be coded by numbers would add
little. This report adds a common categorical language for the interfaces and
proves concrete uniform constructions that this language supports.

\subsection{Main results and their meaning}

The first result is a \emph{category of finite-demand contracts}. An arrow
contains a PR forward transformation and a PR finite set of backward queries,
together with an arithmetic proof that answering those queries suffices.
Composition explicitly substitutes the queries. The categorical equations are
proved from finite-set union and ordinary composition. This has close historical
relatives in Dialectica and Diller--Nahm constructions
\cite{depaiva1991,depaiva2015}; the contribution here is the restricted arithmetic
implementation and its connection to the inspected proof packets, not a claim
to have invented that categorical idea.

The second group of results gives a substantial portion of category theory in
iPRA. For fixed PR presentations, the theory proves elementary diagram
calculations, presentation-wise Yoneda, the unit--counit description of
adjunctions, and an equivalence criterion with explicit splittings. For
uniformly coded finite tables, it additionally constructs functor categories,
finite limits and colimits, finite function objects, bounded choices of
quasi-inverses, and left and right Kan extensions along finite functors.
The ordinary categorical theorems are standard \cite{riehl2016}; the numerical
representations, uniformity boundaries, and induction audits are the point here.

The third result addresses a concrete production gap. Given rational
polynomials and degree caps, a finite matrix test either produces multiplication
identities or a separating linear functional showing that the proposed cap
fails. A successful test yields explicit rational coordinate bounds and a
complete list of rational common zeros. No search over coefficient heights,
unproved Groebner-basis termination, or independent quotient basis is used.
The matrices need not commute. This extends the action-certificate route in
\cite{repo-witnesses}; it does not settle a missing finite-cover or height theorem
for general curves.

Finally, two constructions explain limits that a representation framework
cannot remove. A uniform PR family of decidable equivalence relations, each
with at most two classes, need not admit a computable procedure producing
finite covering lists from the family parameter.
Likewise, a PR binary tree can have an explicitly PR-selected node at every
level but no computable coherent branch. These are precise obstructions to
two tempting uses of categorical language: extracting a finite skeleton merely
from a component count, and replacing a tower of finite observations by a
computably available inverse-limit point.

\subsection{How to read the proof claims}

An \emph{iPRA human proof} in this report specifies numerical operations and
gives an intuitionistic argument whose inductions are quantifier-free or are
explicitly reduced to that case. A \emph{conditional result} retains its
mathematical hypotheses, such as a valid category presentation or a successful
border-closure test. A \emph{Lean target} states what remains to be constructed
in Lean. An executable finite check verifies only the listed finite instances.
These four statuses are not interchangeable.

The archive includes source-level constructions using the real
\texttt{Prf PRAAx} interface for inverse uniqueness, the determining Yoneda
calculation, and logical composition. They have not been compiled in this
environment. The larger arithmetic encodings and main theorem derivations
remain a formalization plan. The user requested at least such an outline;
the paper gives the mathematical arguments independently of those unfinished
formal steps.
